\documentclass[10pt]{article}
 \usepackage[margin=1in]{geometry} 
\usepackage{amsmath,amsthm,amssymb,amsfonts,color, titling}
 \usepackage{listings}

\setlength{\droptitle}{-20mm} 

\usepackage[colorlinks=true,linkcolor=red,urlcolor=blue]{hyperref}

\title{Pre-class assignment \#15}
\author{PHY-905-003\\Computational Astrophysics and
  Astrostatistics\\Spring 2026}
 \date{} % leave blank to have no date

\begin{document}
 
\maketitle

\vspace{-5mm}

\noindent
\textbf{This assignment is due the evening of Monday March 9, 2026.}
 Turn in all materials via GitHub Classroom.

\vspace{5mm}

\noindent
\textbf{Reading/videos:}

\begin{enumerate}

\item Read this article on the
  \href{https://www.bbc.co.uk/bitesize/guides/zws8d2p/revision/1}{basics
    of the CPU}.  Note that there are multiple pages in this article!

  
  
\item Watch
  \href{https://thecrashcourse.com/courses/data-structures-crash-course-computer-science-14/}{this
  video} on data structures. This video is not specific to Python. The first roughly 6 minutes are the most relevant for our next class, but you may find all of it interesting and informative.

  \item Read the accompanying PDF on Lists and Tuples to understand
    how these data structures function under the hood.

\item **Optional, but highly recommended:** Read the
  \href{https://numpy.org/doc/stable/user/absolute_beginners.html}{Beginner's
  Guide to Numpy}

\end{enumerate}

\noindent
\textbf{Your assignment:}  After you read the information and watch
the video above, summarize the key points from them and put any
questions that you might have in the accompanying \texttt{ANSWERS.md}
file

\noindent
\textbf{Handing it in:}
Include your summary and questions in \texttt{ANSWERS.md} and commit it!

\end{document}
